% Observações
% Esse modelo é focado na contestação da reclamação 001/2026 de Pombal, qualquer outra aplicação deve ser adaptada.

\newcommand{\tipoRequerimento}{Contestação de indeferimento de reclamação, faturamento incorreto de iluminação pública}

\section{DOS FATOS E DOS DIREITOS}

\begin{enumerate}[resume]
\item A presente demanda trata de cobrança incorreta e superior ao valor devido na fatura de Iluminação Pública da Unidade Consumidora nº 207926, de responsabilidade do município de Pombal - PB. A desconformidade ocorreu nos meses de março de 2024 e maio de 2025, períodos em que o consumo faturado apresentou-se muito superior ao registrado no Quadro de Iluminação Pública (QIP), resultando em uma diferença de 77.661 kWh cobrados a maior.
\item Em sua resposta administrativa, a concessionária indeferiu o pedido do município sob a alegação de que realizou uma atualização cadastral em seu sistema no dia 08/10/2024 para regularizar os dados do parque de iluminação da cidade. Segundo a empresa, as diferenças nos valores faturados decorrem exclusivamente dessa adequação.
\item No entanto, a justificativa apresentada pela concessionária apresenta uma evidente contradição cronológica. Se a atualização do sistema ocorreu apenas em 08/10/2024, não há fundamentação lógica para que essa alteração justifique ou influencie o valor faturado em março de 2024, sete meses antes da modificação do cadastro. Uma atualização posterior não possui o condão de validar cobranças feitas em períodos passados.
\item Além disso, o argumento também se mostra inconsistente em relação ao mês de maio de 2025. Mesmo após a suposta correção do sistema em outubro de 2024, a fatura de maio de 2025 registrou um consumo de 115.873 kWh, enquanto o QIP daquele mesmo mês indicava apenas 59.995 kWh. Constata-se, portanto, que a concessionária continuou a realizar a cobrança em desacordo com os parâmetros vigentes, faturando quase o dobro do valor correto.
\item Como o erro no faturamento foi de responsabilidade exclusiva da concessionária e a justificativa apresentada não possui consistência cronológica ou matemática, fica demonstrada a irregularidade da cobrança. Sendo assim, o município possui o direito à restituição em dobro dos valores pagos a maior, nos termos estabelecidos pela Resolução Normativa nº 1.000/2021 da ANEEL.
\end{enumerate}

\section{DOS PEDIDOS}

\begin{enumerate}[resume]
\item Diante do exposto, requer-se:
    \begin{enumerate}
    \item O recebimento e provimento desta contestação para reformar a decisão da concessionária, reconhecendo-se a incorreção das cobranças efetuadas nos meses de março de 2024 e maio de 2025;
    \item A determinação para que a concessionária realize a devolução em dobro do valor equivalente aos 77.661 kWh cobrados a maior (totalizando 155.322 kWh), com a devida atualização monetária pelo IPCA e aplicação de juros, conforme as normas da ANEEL;
    \item Que a referida restituição seja efetuada por meio de depósito bancário diretamente na conta corrente da Prefeitura Municipal de Pombal - PB;
    \item Que a concessionária seja compelida a disponibilizar o histórico completo de consumo dos últimos 115 meses da referida unidade consumidora ou, alternativamente, ao menos uma fatura correspondente ao mês de dezembro de cada ano desse período.
    \end{enumerate}
\end{enumerate}